\documentclass[11pt,a4paper]{article}

\usepackage[utf8]{inputenc}
\usepackage[T1]{fontenc}
\usepackage{lmodern}
\usepackage{textcomp}
\providecommand{\euro}{\texteuro}
\usepackage[french]{babel}
\usepackage[margin=1.9cm]{geometry}
\usepackage{microtype}
\usepackage{booktabs}
\usepackage{array}
\usepackage{amsmath,amssymb}
\usepackage[table]{xcolor}
\usepackage[most]{tcolorbox}

\definecolor{navy}{HTML}{1F3864}
\definecolor{lightblue}{HTML}{D6E0F0}
\definecolor{accent}{HTML}{C0491F}
\definecolor{softgreen}{HTML}{E2EFDA}

\newtcolorbox{cle}[1][]{colback=lightblue!40,colframe=navy,boxrule=0.6pt,
  arc=2pt,left=8pt,right=8pt,top=6pt,bottom=6pt,#1}

\newcommand{\oui}{\textcolor{navy}{$\checkmark$}}
\newcommand{\non}{\textcolor{accent}{$\times$}}
\newcommand{\cascell}{\cellcolor{softgreen}}

\setlength{\parindent}{0pt}
\setlength{\parskip}{4pt}

\title{\vspace{-1.3cm}\color{navy}\bfseries Positionnement et revue de littérature\\[2pt]
  \large La dépendance à l'ordre, un mécanisme absent des approches existantes}
\author{Kélian Kaddouri}
\date{\today}

\begin{document}
\maketitle
\vspace{-0.7cm}

\begin{cle}
\textbf{La contribution en une phrase.} Là où la littérature cyber actuarielle modélise soit
la \emph{co-occurrence} (copules, chocs communs), soit la \emph{contagion entre entités}
(réseaux, épidémiologie), soit la \emph{dynamique temporelle} (Hawkes), ce mémoire introduit
une \textbf{dépendance dirigée et à l'ordre entre les domaines de contrôle réglementaires}
(les piliers DORA) au sein d'une même entité, et la traduit en SCR. Sa signature est que
la criticité dépend de l'\textbf{ordre} de la chaîne et non du seul ensemble des piliers
touchés, ce qu'aucune des approches recensées ne peut reproduire.
\end{cle}

\section*{1. Le mécanisme, et sa signature}

La cascade dirigée pose une matrice $W$ de contagion entre piliers, où $W_{jk}$ mesure la
propension d'une défaillance de $k$ à entraîner celle de $j$. Trois propriétés la
distinguent : elle est \textbf{dirigée} ($W_{jk}\neq W_{kj}$), elle autorise les
\textbf{cycles} (P1$\to$P2$\to\cdots\to$P1), stabilisés par la normalisation de Leontief
($\rho(W)\le g<1$, forme réduite $(I-W)^{-1}$), et surtout elle est \textbf{à l'ordre} : la
criticité d'une chaîne n'est pas une fonction de l'\emph{ensemble} des piliers touchés, elle
dépend de l'\emph{ordre} dans lequel ils le sont. Ce n'est pas une intuition : sur les 26
sous-ensembles de taille au moins deux, \textbf{22 changent de criticité selon l'ordre},
avec un écart atteignant \textbf{6 points sur une échelle de 10} (la chaîne P1$\to$P2 vaut
8, la chaîne P2$\to$P1 vaut 2). Or toute dépendance symétrique et tout score additif sont,
par construction, invariants par permutation : \textbf{aucun ne peut reproduire cela}.

\begin{cle}
\textbf{Remarque de méthode, et ce qui n'est \emph{pas} revendiqué.} Une propriété plus
forte, la \emph{non-transitivité} de la dominance entre piliers, a été testée dans trois
formulations : la transitivité des corrélations qu'impose un modèle à facteur unique, la
positivité de la dépendance symétrisée, et l'existence de cycles de dominance par paire.
\textbf{Aucune ne se vérifie} sur ce modèle. Elle n'est donc pas revendiquée. La
dépendance à l'ordre, elle, se vérifie, et elle exclut exactement les mêmes concurrents.
Signaler l'énoncé que l'on a testé et abandonné vaut mieux que de le laisser à la
sagacité du lecteur.
\end{cle}

\section*{2. Positionnement : pourquoi le mécanisme est irréductible}

\begin{center}\footnotesize
\renewcommand{\arraystretch}{1.35}
\begin{tabular}{@{}p{3.2cm} c c c c p{3.1cm} p{3.0cm}@{}}
\toprule
\textbf{Approche} & \textbf{Dir.} & \textbf{Ordre} & \textbf{Cycles} & \textbf{Propag.} &
\textbf{Donnée requise} & \textbf{Objet modélisé} \\
\midrule
Matrices de risque / AMDEC & \non & \non & \non & \non & dires d'expert & risques listés (proba $\times$ gravité) \\
Copule (Gumbel, gaussienne) & \non & \non & \non & \non & co-mouvement observé & briques de perte \\
Réseau / épidémiologie (SIR) & \oui & \non & \oui & \oui & structure réseau $+$ taux d'infection & entités d'un portefeuille \\
Hawkes (multivarié) & \oui & \non & \oui & \oui & horodatage fin des arrivées & types d'attaque, dans le temps \\
Réseau bayésien & \oui & \non & \non & \oui & tables de proba conditionnelle & variables d'un graphe acyclique \\
\cascell\textbf{Cascade dirigée (ici)} & \cascell\oui & \cascell\oui & \cascell\oui & \cascell\oui & \cascell jugement d'expert (élicité) & \cascell\textbf{domaines de contrôle (piliers) d'une entité} \\
\bottomrule
\end{tabular}
\end{center}

\vspace{2pt}
Point par point :
\begin{itemize}\itemsep2pt
\item \textbf{Matrices de risque / AMDEC} (le standard praticien, RPN $=$ gravité $\times$
      occurrence $\times$ détection) : listent et notent des risques mais n'ont \emph{aucune}
      propagation ni direction. La cascade ajoute exactement la couche qui leur manque.
\item \textbf{Copules} (Böhme et Kataria 2006 ; Herath et Herath 2011) : capturent la
      dépendance de queue entre pertes, mais \textbf{symétrique} : elles ne distinguent pas
      $k\to j$ de $j\to k$. Notre benchmark (script 23) le chiffre : une copule ne peut pas
      exprimer la propagation (interaction moyenne $+0$ par construction contre $+560$ M\euro{}
      pour la cascade), et sa priorité de remédiation suit les \emph{réceptacles}, la cascade
      les \emph{sources}.
\item \textbf{Modèles réseau et épidémiologiques} (Hillairet et Lopez 2021, 2022 ; Xu et Hua
      2017 ; Baldwin et al.\ 2017) : ils modélisent la vraie contagion, mais \textbf{entre
      entités} d'un portefeuille, et \textbf{exigent la structure du réseau} et des taux
      d'infection, données indisponibles ici. Nous \emph{empruntons} à Hillairet et Lopez la
      logique d'états de conformité, mais l'appliquons \textbf{entre piliers} d'une entité,
      pas entre assurés.
\item \textbf{Hawkes} (Bessy-Roland, Boumezoued et Hillairet 2021 ; Boumezoued, Cherkaoui et
      Hillairet 2023) : modélisent l'auto-excitation de la \textbf{fréquence dans le temps},
      entre \emph{types d'attaque}, pas la structure inter-piliers. Nous montrons (scripts
      08c, 08h) que sur la fréquence d'entrée, l'excitation est intra-journalière et exogène,
      et nous ne l'ajoutons pas ; ce constat rejoint leur raffinement two-phase de 2023.
\item \textbf{Réseaux bayésiens} : peuvent porter la direction, mais supposent un graphe
      \textbf{acyclique} (pas de cycle de contagion) et des tables de probabilités
      conditionnelles qui se heurtent au \textbf{même mur d'identification} que $W$ (chapitre
      identifiabilité). La forme de Leontief, elle, gère les cycles.
\end{itemize}

\section*{3. Revue de littérature, par thème}

\textbf{Dépendance et accumulation.} La modélisation actuarielle du cyber a d'abord traité la
\emph{corrélation} entre risques par copules (Böhme et Kataria 2006 ; Herath et Herath 2011)
et par chocs communs. Ces approches capturent l'accumulation (que nos données confirment,
concentration journalière de Gini $0{,}69$, cf.\ cas d'usage P4) mais restent symétriques.

\textbf{Contagion structurée.} Une seconde vague introduit la propagation : Hillairet et Lopez
(2021) couplent des processus de comptage à un modèle épidémiologique compartimental (type
SIR) pour la contagion dans un portefeuille ; Hillairet, Lopez, d'Oultremont et Spoorenberg
(2022) ajoutent une structure de réseau ; Xu et Hua (2017) et Baldwin et al.\ (2017)
modélisent la contagion entre cibles. Toutes exigent la donnée de réseau ou d'infection.

\textbf{Dynamique temporelle.} La fréquence auto-excitée est traitée par les processus de
Hawkes : Bessy-Roland, Boumezoued et Hillairet (2021) sur données PRC (dont notre source
dérive), puis Boumezoued, Cherkaoui et Hillairet (2023) avec un terme exogène ; Hillairet,
Réveillac et Rosenbaum (2021) en développent la théorie pour le pricing de dérivés ; Peng et
al.\ (2016) sur des taux d'attaque horodatés. Notre travail se situe dans cette lignée mais
en \emph{sort} sur la fréquence d'entrée, pour les raisons données au chapitre fréquence.

\textbf{Fréquence, sévérité, queue.} Sur les marges, Edwards et al.\ (2016) ajustent des lois
discrètes (binomiale négative) au nombre journalier de brèches ; Lu, Zhang et Zhu (2024)
séparent fréquence et nombre de brèches par événement (structure de paquet qui soutient notre
Poisson composé) ; Farkas, Lopez et Thomas (2019) tarifent par arbres de régression GPD. Nous
en retenons la théorie des valeurs extrêmes (indice de queue $\xi$) et la structure de paquet.

\textbf{Cadre structurel.} Le squelette quantitatif emprunte au modèle à facteur de Vasicek
(risque de crédit) l'idée d'un stress latent à composante systémique, et à l'analyse
entrées-sorties de Leontief la normalisation qui garantit la stabilité d'une contagion
\emph{avec cycles}. La nouveauté est de faire porter cette mécanique par une matrice
\textbf{dirigée} entre domaines de contrôle réglementaires.

\section*{4. Où se situe ce mémoire}

Le mémoire \textbf{emprunte} ce qui est établi et ancré sur données (la queue de sévérité par
EVT, la structure de paquet de la fréquence, la logique d'états de conformité de Hillairet et
Lopez) et \textbf{ajoute} deux choses absentes de la littérature : (i)~une dépendance
\textbf{dirigée et à l'ordre} entre les piliers DORA, traduite en
capital ; (ii)~une \textbf{frontière d'identifiabilité} explicite, muée en spécification de
reporting. Aucune des références recensées ne modélise une dépendance dirigée entre domaines
de contrôle au sein d'une entité, ni ne la relie au SCR.

\begin{cle}
\textbf{La revendication de nouveauté, défendable.} \og Je n'invente pas une nouvelle loi de
sévérité ni un nouveau processus de fréquence : sur ces marges, je m'appuie sur la
littérature établie et je la calibre. Ma contribution est le \textbf{mécanisme de dépendance}
lui-même, dirigé et à l'ordre entre piliers réglementaires, que ni les copules (symétriques),
ni les modèles de réseau (qui exigent une donnée absente), ni les processus de Hawkes (qui
vivent dans le temps) ne capturent, et que je traduis en capital tout en démontrant sa limite
d'identification. \fg{} C'est une nouveauté \emph{maîtrisée} : neuve sur le mécanisme, sobre
sur le reste, honnête sur les limites.
\end{cle}

\section*{Références (à compléter en style institut des actuaires)}
\footnotesize
Baldwin, Gheyas, Ioannidis, Pym, Williams (2017), \emph{Contagion in cyber security attacks},
J.\ Operational Research Society 68(7). \quad
Bessy-Roland, Boumezoued, Hillairet (2021), \emph{Multivariate Hawkes process for cyber
insurance}, Annals of Actuarial Science 15(1). \quad
Böhme, Kataria (2006), \emph{Models and measures for correlation in cyber-insurance}, WEIS. \quad
Boumezoued, Cherkaoui, Hillairet (2023), \emph{Cyber risk modeling using a two-phase Hawkes
process with external excitation}, arXiv:2311.15701. \quad
Edwards, Hutchings, Forrest (2016), \emph{Hype and heavy tails: a closer look at data
breaches}, J.\ Cybersecurity 2. \quad
Farkas, Lopez, Thomas (2019), \emph{Cyber claim analysis through generalized Pareto regression
trees}. \quad
Herath, Herath (2011), \emph{Copula based actuarial model for pricing cyber-insurance
policies}, Insurance Markets and Companies 2. \quad
Hillairet, Lopez (2021), \emph{Propagation of cyber incidents in an insurance portfolio},
Scandinavian Actuarial Journal 2021(8). \quad
Hillairet, Lopez, d'Oultremont, Spoorenberg (2022), \emph{Cyber-contagion model with network
structure applied to insurance}, Insurance: Mathematics and Economics 107. \quad
Hillairet, Réveillac, Rosenbaum (2021), \emph{An expansion formula for Hawkes processes\dots},
arXiv:2104.01579 (SPA 2023). \quad
Lu, Zhang, Zhu (2024), \emph{Cyber risk modeling: a discrete multivariate count process
approach}, Scandinavian Actuarial Journal 2024(6). \quad
Peng, Xu, Xu, Hu (2016), \emph{Modeling and predicting extreme cyber attack rates via marked
point processes}, J.\ Applied Statistics 44(14). \quad
Xu, Hua (2017), \emph{Cybersecurity insurance: modeling and pricing}, SOA. \quad
Vasicek (2002), \emph{The distribution of loan portfolio value}, Risk. \quad
Leontief (1936), \emph{Quantitative input and output relations\dots}.
\end{document}
